% chapters/conclusion.tex — Conclusion chapter (narrative text only)
\chapter{Conclusion}
\label{chap:conclusion}

This survey has examined five core capabilities of reasoning-centric agentic
\LLM{}s: planning, memory, retrieval, tool use, and multimodal reasoning.
Across all five, a consistent theme emerged: agents that decompose tasks
into bounded sub-problems, maintain explicit state, and ground decisions in
retrieved evidence substantially outperform agents that rely on a single,
monolithic forward pass~\cite{wei2026agenticreasoning,ke2025reasoningfrontiers}.
\index{agentic reasoning}

\section{Key Findings}
\label{sec:conclusion-findings}

\begin{enumerate}
  \item \textbf{Planning:} Search-augmented frontier \LLM{}s approach or
        match state-of-the-art classical planners on standard
        benchmarks~\cite{correa2025planningperformance}, suggesting that
        the gap between neural and symbolic planning is narrowing rapidly.

  \item \textbf{Memory:} Lightweight cognitive architectures
        (\cite{huang2025licomemory,chen2025telemem}) provide an effective
        middle ground between context-window memory and external databases,
        enabling long-horizon tasks without prohibitive storage overhead.

  \item \textbf{Retrieval:} Proactive retrieval strategies reduce latency
        and improve grounding accuracy, particularly in high-stakes
        domains~\cite{wang2025proactiveretrievalmedical}.

  \item \textbf{Tool use:} Bounded tool-use policies combined with
        code-execution loops enable \LLM{}s to solve tasks that require
        exact computation~\cite{liu2025agenticmath}.

  \item \textbf{Multimodal:} Dynamic modality-fusion architectures
        consistently outperform static alternatives, pointing toward
        unified agentic reasoners that treat all input channels
        symmetrically~\cite{yao2025multimodalsurvey}.
\end{enumerate}

\section{Future Directions}
\label{sec:conclusion-future}

Open problems span all five capability axes.
For planning, learning to set task horizons adaptively remains unsolved.
For memory, efficient forgetting — retaining task-relevant facts while
discarding irrelevant detail — is an active research question.
For retrieval, the alignment between retrieval relevance and downstream
task reward deserves more theoretical treatment.
For tool use, compositional policies that chain multiple tools in a single
reasoning step are in their infancy.
For multimodal reasoning, benchmarks that isolate cross-modal inference
from language understanding are needed.

Progress on each of these fronts will bring agentic \LLM{}s closer to the
vision articulated by \cite{wei2026agenticreasoning}: systems that reason,
plan, and act with the generality and robustness required for open-ended,
real-world deployment.
